% =============================================================================
% 6. RELATED WORK
% =============================================================================
\section{Related Work}
\label{sec:related}

\paragraph{Mid-training and continued pretraining.}
A growing line of work has identified a distinct training stage between
large-scale pretraining and task-specific post-training, in which a model is
further trained on a smaller but more carefully curated corpus to install
inductive biases that are difficult to acquire from either generic web text
or downstream fine-tuning data. MiniCPM~\citep{minicpm} and
OLMo~\citep{olmo} explicitly schedule such a stage to upweight high-quality
or domain-relevant data near the end of pretraining, and
Code-Llama~\citep{codellama} long-context adaptation can be read as a
mid-training stage targeting context-length generalisation. We adopt the
same staging philosophy but apply it to a different end: rather than
optimising for general data quality or context length, we use mid-training
as the natural place to install an \textit{agent-oriented} structural prior
through function-aware FIM. Positioning the objective at this stage
concentrates its signal immediately before agentic post-training rather
than diluting it across trillions of pretraining tokens.

\paragraph{Fill-in-the-middle objectives.}
The fill-in-the-middle objective~\citep{bavarian2022fim} has become a
standard ingredient in modern code LLM pretraining, with random-span
variants mixed into the pretraining corpora of
Code-Llama~\citep{codellama}, StarCoder~\citep{starcoder},
Qwen-Coder~\citep{qwencoder}, and DeepSeek-Coder~\citep{deepseekcoder}.
Our work builds directly on this body of evidence that bidirectional
infilling is a useful complement to left-to-right modelling, and we
explicitly view our recipe as an \textit{extension and specialisation} of
prior FIM rather than a replacement. Three differences mark our setting.
First, masking targets are chosen at \textit{function granularity} via
program dependency graph analysis and a complexity--inferability double
criterion, so that the masked region corresponds to a unit of reasoning
rather than a syntactically arbitrary span. Second, we embed an
\textit{explicit chain-of-thought} inside the FIM middle so that the
model is supervised to produce a rationale before the code, mirroring the
think-then-act structure of an agent step. Third, the objective is applied
in a \textit{dedicated mid-training stage} rather than dissolved into
pretraining, so that the structural prior reaches post-training in
concentrated form. None of the three differences is in tension with prior
random-span FIM; together, they specialise the same family of objectives
to coding-agent capability.

\paragraph{Coding agent foundation models.}
Recent progress on coding agents has been driven by benchmarks such as
SWE-Bench~\citep{swebench}, scaffolds such as
SWE-agent~\citep{sweagent} and OpenHands~\citep{openhands}, and
trajectory-centric post-training pipelines such as
SWE-Gym~\citep{swegym}, r2e-gym~\citep{r2egym}, and
SWE-Smith~\citep{swesmith}. These pipelines synthesise or curate agent
trajectories that imitate human or LLM behaviour on issue-resolution
tasks, and they have repeatedly demonstrated that high-quality trajectory
data is sufficient to push end-task numbers. Our work is complementary
rather than competing: the post-training pipelines we use in our main
results are r2e-gym and SWE-Smith, applied without modification on top of
a mid-trained checkpoint. The novelty lies one stage earlier. Instead of
producing more synthetic trajectories, we extract a self-supervised
signal from the structure already present in ordinary source code, and
we show that this signal both raises in-domain agent metrics and
mitigates the off-domain capability erosion that trajectory-only
post-training otherwise inflicts.

\paragraph{Distillation from frontier models.}
Our default recipe uses Gemini-3-Preview to generate the chain-of-thought
rationales embedded inside each FIM middle span, placing it within the
broader tradition of distilling capability from a strong teacher into a
smaller student~\citep{orca, selfinstruct, distillwhisper}. We
acknowledge this dependency upfront because it is essential for an honest
reading of our results. The role of the rationale here, however, is not
to transfer the teacher's generic reasoning ability but to provide a
structured intermediate supervision aligned with the FIM objective.
We isolate the contribution of FIM structure from that of CoT
distillation through a controlled ablation in
Section~\ref{sec:exp:ablation}: removing the rationale entirely still
beats the no-mid-training baseline, and replacing the frontier teacher
with self-generated rationales recovers a sizeable fraction of the gain.
We therefore do not view the recipe as fundamentally distillation-driven,
even though the best configuration does benefit from teacher rationales.
